\uuid{l19f}
\titre{ Interprétation des courbes de coût lors de l'entraînement }
\niveau{L3}
\module{Analyse de données}
\chapitre{Réseaux de neurones}
\sousChapitre{Descente de gradient, diagnostic d'entraînement}
\theme{Courbe de coût, learning rate, divergence, neurone mort, initialisation}
\auteur{Maxime Nguyen}
\datecreate{2026-03-19}
\organisation{AMSCC}
\difficulte{3}

\contenu{
\texte{
  On entraîne un modèle de régression linéaire par descente de gradient et on observe trois allures possibles de la courbe de coût :
  \begin{itemize}
    \item \textbf{Courbe A} : descend régulièrement et se stabilise.
    \item \textbf{Courbe B} : descend puis remonte en explosant.
    \item \textbf{Courbe C} : reste parfaitement plate dès la première itération.
  \end{itemize}
  Parmi les causes suivantes, associer la plus probable à chaque courbe :
  \begin{enumerate}[(1)]
    \item \texttt{lr = 10} au lieu de \texttt{lr = 0.01}
    \item \texttt{w -= lr * dw} remplacé par \texttt{w += lr * dw}
    \item $w$ et $b$ initialisés à zéro sur un réseau à une couche cachée
    \item Tout est correct
  \end{enumerate}
}
\begin{enumerate}

\item
  \question{Associer la cause la plus probable à chacune des trois courbes. Justifier.}
  \reponse{
    \begin{itemize}
      \item \textbf{Courbe A} $\to$ cause (4) : tout est correct, l'entraînement converge normalement.
      \item \textbf{Courbe B} $\to$ cause (1) ou (2) : un learning rate trop élevé ($\alpha = 10$) provoque
            des mises à jour trop grandes et fait diverger le coût ; de même, \texttt{w += lr * dw}
            effectue une montée de gradient, ce qui fait croître $J$.
      \item \textbf{Courbe C} $\to$ cause (3) : initialiser tous les poids à zéro sur un réseau à couche cachée
            entraîne une symétrie parfaite — tous les neurones sont identiques, les gradients s'annulent
            ou se répètent, et le coût ne bouge pas (ou très peu).
    \end{itemize}
  }

\end{enumerate}
}
